% Czech and fonts
\documentclass[11pt, a4paper]{article}
\usepackage[utf8]{inputenc} 
\usepackage[T1]{fontenc} 
\usepackage[czech]{babel}
\usepackage{lmodern}

% Margins and colors
\usepackage[top=2.5cm, bottom=2.5cm, left=3cm, right=2.5cm]{geometry}
\usepackage{graphicx} 
\usepackage{xcolor} 

% Formating of pseudocode 
\usepackage{listings} 
\lstset{ backgroundcolor=\color{white}, basicstyle=\ttfamily\small, breaklines=true, frame=none, xleftmargin=1cm, xrightmargin=0.5cm, aboveskip=4pt, belowskip=4pt, }

% Set algorithms 
\usepackage{algorithm} 
\usepackage{algpseudocode} 
\usepackage{float} % Nutne pro parametr [H] u obrazku
\floatname{algorithm}{Algorithm} 
\renewcommand{\algorithmicrequire}{\textbf{Input:}} 
\renewcommand{\algorithmicensure}{\textbf{Output:}}

% Hyperlinks 
\usepackage[hidelinks]{hyperref}

% Set header and footer 
\usepackage{fancyhdr} 
\pagestyle{fancy} 
\fancyhf{} 
\renewcommand{\headrulewidth}{0.4pt} 
\renewcommand{\footrulewidth}{0pt} 
\fancyhead[C]{\small\itshape Algoritmy počítačové kartografie \quad Digitální model terénu} 
\fancyfoot[C]{\thepage}

% Set spacing
\usepackage{setspace} 
\setstretch{1.5}

\usepackage{amsmath}

\begin{document}

% Title page
\begin{titlepage}
    \centering

    {\large Přírodovědecká fakulta}\\[0.3cm]
    {\large Univerzita Karlova}\\[3cm]

    \includegraphics[width=0.425\textwidth]{image.png}\\[4cm]

    {\large Algoritmy počítačové kartografie}\\[0.5cm]
    {\LARGE\bfseries Digitální model terénu}\\[4cm]

    {\large Kristýna Antošová, Vít Kadlec}\\[0.5cm]
    {\normalsize 1.N-GKDPZ}\\[0.3cm]
    {\normalsize Praha 2026}

    \vfill
\end{titlepage}

% Assignment 
\section{Zadání}
\medskip 
\begin{figure}[H] 
\centering 
\includegraphics[width=\textwidth]{zadani.png} 
\end{figure}
 
\begin{figure}[H] 
\centering 
\includegraphics[width=\textwidth]{Hodnoceni.png} 
\end{figure}

\noindent V rámci zadání byly zpracovány první dvě základní úlohy: konstrukce Delaunay triangulace s polyedrickým modelem terénu a konstrukce vrstevnic doplněná o analýzu sklonu a expozice. Zpracována byla navíc úloha: Výběr barevných stupnic při vizualizaci sklonu a expozice.
\newpage

% Technical report 
\section{Popis a rozbor problému}

\subsection{Digitální model terénu}
Digitální model terénu (DMT) je digitální reprezentace zemského povrchu tvořená množinou bodů $P = \{p_1, \ldots, p_n\}$, kde $p_i = (x_i, y_i, z_i)$. Nad touto množinou se konstruuje polyedrický model složený z rovinných trojúhelníkových plošek (TIN -- Triangulated Irregular Network). Výhodou TIN oproti rastrové reprezentaci je, že zachovává přesnou polohu vstupních bodů a přizpůsobuje hustotu ploch členitosti terénu.

Nad takto vytvořeným modelem lze provádět celou řadu analytických úloh -- tvorbu vrstevnic, analýzu sklonu, analýzu expozice (orientace ke světovým stranám), výpočet objemů zemních prací, viditelnosti a další. Základním stavebním kamenem je však triangulace samotná, jejíž kvalita zásadně ovlivňuje všechny následné výstupy.

\subsection{Delaunay triangulace}
Delaunay triangulace je speciální typ triangulace, která maximalizuje minimální úhel ve všech trojúhelnících. Díky této vlastnosti poskytuje vizuálně i numericky nejkvalitnější výsledky -- minimalizuje výskyt dlouhých úzkých trojúhelníků, které jsou pro interpolaci terénu nevhodné.

Definičním kritériem je tzv. \textbf{kritérium prázdné kružnice}: kružnice opsaná libovolnému trojúhelníku triangulace neobsahuje žádný další bod vstupní množiny. Ekvivalentním kritériem, které je snadněji algoritmicky ověřitelné, je maximalizace úhlu nad protilehlou hranou: pro hranu $(p_1, p_2)$ hledáme takový bod $p_{DT}$, pro který platí:
\begin{equation}
p_{DT} = \arg \max_{\forall p_i \in P} \angle(p_1, p_i, p_2).
\end{equation}

\subsubsection{Inkrementální konstrukce}
Mezi hlavní metody konstrukce Delaunay triangulace patří \textit{Inkrementální konstrukce}, \textit{Divide and Conquer} a \textit{Flip algorithm}. V této úloze byla implementována \textbf{inkrementální konstrukce} s využitím seznamu aktivních hran (AEL -- Active Edge List).

Algoritmus začíná nalezením pivotu $q$ (bod s minimální $y$-souřadnicí) a jeho nejbližšího souseda $q_n$. Hrana $(q, q_n)$ leží na konvexní obálce a tvoří výchozí hranu triangulace. K této hraně je v levé polorovině hledán bod $p_{DT}$ splňující Delaunayovo kritérium. Tím vzniká první trojúhelník. Nové hrany jsou přidány do AEL; pokud je některá hrana v AEL nalezena s opačnou orientací, znamená to, že byla zpracována z obou stran, a je z AEL odstraněna. Proces pokračuje, dokud AEL neobsahuje žádné hrany.

\begin{figure}[H] 
\centering 
\includegraphics[width=0.6\textwidth]{delaunay.png} 
\caption{\textit{Princip inkrementální konstrukce Delaunay triangulace (zdroj: Bayer 2026)}}
\end{figure}

\subsection{Konstrukce vrstevnic}
Vrstevnice je čára spojující body o stejné nadmořské výšce. Matematicky se jedná o průsečnici roviny $z = z_k$ (konstantní výška) s plochou DMT. Nad TIN se vrstevnice konstruují pomocí lineární interpolace po jednotlivých trojúhelnících.

Pro každý trojúhelník s vrcholy $p_1, p_2, p_3$ a zvolenou výšku $z$ se ověří, zda daná rovina trojúhelník protíná. Označme výškové rozdíly:
\begin{equation}
\Delta z_i = z - z_i, \quad i \in \{1, 2, 3\}.
\end{equation}
Rovina protíná trojúhelník, pokud mají výškové rozdíly pro dvě různé hrany opačné znaménko. Průsečík $b$ hrany $(p_1, p_2)$ s rovinou je dán lineární interpolací:
\begin{equation}
x_b = \frac{x_2 - x_1}{z_2 - z_1}(z - z_1) + x_1, \quad y_b = \frac{y_2 - y_1}{z_2 - z_1}(z - z_1) + y_1.
\end{equation}
Obdobně se vypočte druhý průsečík, a spojnice mezi nimi tvoří jeden segment vrstevnice v daném trojúhelníku.

V aplikaci se rozsah nadmořských výšek $z_{min}$ a $z_{max}$ automaticky určuje z minimální a maximální hodnoty $z$ v načtené množině bodů. Krok $dz$ mezi jednotlivými vrstevnicemi je volen tak, aby napříč celým výškovým rozsahem vzniklo přibližně 20 vrstevnic, což zajišťuje čitelnou vizualizaci nezávisle na konkrétních hodnotách vstupních dat.

Pro dodržení kartografických zásad je implementováno rozlišení \textbf{zvýrazněných (hlavních) vrstevnic}. Hlavní vrstevnice jsou vykreslovány každou pátou úroveň (tj. s intervalem $5 \cdot dz$) silnější tmavší hnědou čarou, zatímco ostatní (vedlejší) vrstevnice jsou vykreslovány tenčí světlejší hnědou čarou. Takové odlišení usnadňuje orientaci v mapě a rychlou identifikaci hlavních výškových úrovní.

\subsection{Analýza sklonu}
Sklon terénu je úhel $\varphi$ mezi rovinou trojúhelníku $\rho$ a horizontální rovinou $\pi$. Je to jedna z nejdůležitějších analytických úloh nad DMT s využitím v hydrologii, zemědělství, návrhu komunikací či posuzování stability svahů.

Rovina $\rho$ procházející třemi body $p_1, p_2, p_3$ je určena normálovým vektorem:
\begin{equation}
\vec{n}_1 = (a, b, c) = \vec{p_1 p_2} \times \vec{p_1 p_3}.
\end{equation}
Normálový vektor horizontální roviny je $\vec{n}_2 = (0, 0, 1)$. Úhel $\varphi$ mezi rovinami je roven úhlu mezi jejich normálami:
\begin{equation}
\varphi = \arccos \frac{\vec{n}_1 \cdot \vec{n}_2}{|\vec{n}_1| \cdot |\vec{n}_2|} = \arccos \frac{c}{\sqrt{a^2 + b^2 + c^2}}.
\end{equation}
Hodnota $\varphi = 0$ odpovídá vodorovné ploše, $\varphi = \pi/2$ svislé stěně. Výsledek se vizualizuje mapováním sklonu do stupnice šedi -- čím strmější svah, tím tmavší barva.

\begin{figure}[H] 
\centering 
\includegraphics[width=0.65\textwidth]{sklon.png} 
\caption{\textit{Geometrická interpretace sklonu terénu (zdroj: Bayer 2026)}}
\end{figure}

\subsection{Analýza expozice}
Expozice (orientace) terénu je definována jako azimut průmětu gradientu $\nabla \rho$ do roviny $xy$. Udává, ke které světové straně je daný svah orientován. Pro severní svah je expozice $0^{\circ}$, východní $90^{\circ}$, jižní $180^{\circ}$ a západní $270^{\circ}$.

Průmět gradientu do roviny $xy$ má tvar:
\begin{equation}
\vec{v} = \left(\frac{\partial \rho}{\partial x}(x_0), \frac{\partial \rho}{\partial y}(y_0), 0\right) = (a, b, 0).
\end{equation}
Azimut $A$ měřený od osy $y$ (sever) je dán vztahem $A = \arctan(a/b)$. V implementaci se pro správné určení kvadrantu a normalizaci výsledku do intervalu $[0, 2\pi)$ používá funkce \texttt{atan2}. Protože je při zobrazení na obrazovce $y$-ová osa převrácená (roste směrem dolů, nikoliv nahoru), je nutné pro zachování geografické orientace negovat $y$-ovou složku vektoru:
\begin{equation}
A = \text{atan2}(a, -b).
\end{equation}
Tato korekce zajišťuje, že severní svahy skutečně odpovídají severu na obrazovce. Výstup je vizualizován pomocí diskrétní barevné palety rozdělující plný azimut na 8 hlavních světových stran.

Expozice terénu má význam v zemědělství (orientace ovlivňuje množství dopadajícího slunečního záření, a tím i růst plodin), stavebnictví, hydrologii i energetice (návrh solárních elektráren).

\begin{figure}[H] 
\centering 
\includegraphics[width=0.55\textwidth]{expozice.png} 
\caption{\textit{Geometrická interpretace expozice terénu (zdroj: Bayer 2026)}}
\end{figure}

\newpage
\section{Popisy algoritmů formálním jazykem}

\subsection{Inkrementální konstrukce Delaunay triangulace}
\begin{algorithm}[H] \caption{Incremental Delaunay Triangulation}
\begin{algorithmic}[1]
\Require body $P$
\Ensure Delaunay triangulace $DT$
\State inicializuj $DT \gets \emptyset$, $AEL \gets \emptyset$
\State nalezni pivot $q$ s minimální $y$-souřadnicí
\State nalezni bod $q_n$ nejbližší k $q$
\State vytvoř hranu $e = (q, q_n)$ a $e_s = (q_n, q)$, přidej obě do $AEL$
\While{$AEL$ není prázdný}
    \State odeber hranu $e_1$ z $AEL$
    \State přehoď orientaci: $e_{1s} = e_1.switchOrientation()$
    \State nalezni Delaunay bod $p_{DT}$ v levé polorovině od $e_{1s}$
    \If{$p_{DT}$ neexistuje} \textbf{continue} \EndIf
    \State vytvoř nové hrany $e_2 = (e_{1s}.end, p_{DT})$, $e_3 = (p_{DT}, e_{1s}.start)$
    \State přidej $e_{1s}, e_2, e_3$ do $DT$
    \State aktualizuj $AEL$ pro $e_2$ a $e_3$
\EndWhile
\end{algorithmic}
\end{algorithm}

\subsection{Nalezení Delaunay bodu}
\begin{algorithm}[H] \caption{Find Delaunay Point}
\begin{algorithmic}[1]
\Require body $P$, hrana $(p_1, p_2)$
\Ensure bod $p_{DT}$
\State $p_{DT} \gets null$, $\varphi_{max} \gets 0$
\For{každý bod $p_i \in P$, $p_i \neq p_1, p_2$}
    \If{$p_i$ leží v levé polorovině od $(p_1, p_2)$}
        \State $\varphi \gets \angle(p_1, p_i, p_2)$
        \If{$\varphi > \varphi_{max}$}
            \State $\varphi_{max} \gets \varphi$, $p_{DT} \gets p_i$
        \EndIf
    \EndIf
\EndFor
\end{algorithmic}
\end{algorithm}

\subsection{Aktualizace seznamu aktivních hran}
\begin{algorithm}[H] \caption{Update AEL}
\begin{algorithmic}[1]
\Require hrana $e$, seznam $AEL$
\State $e_s \gets e.switchOrientation()$
\If{$e_s \in AEL$}
    \State odeber $e_s$ z $AEL$
\Else
    \State přidej $e$ do $AEL$
\EndIf
\end{algorithmic}
\end{algorithm}

\subsection{Konstrukce vrstevnic}
\begin{algorithm}[H] \caption{Create Contour Lines}
\begin{algorithmic}[1]
\Require triangulace $DT$, $z_{min}$, $z_{max}$, krok $dz$
\Ensure vrstevnice $contours$
\For{$z$ od $z_{min}$ do $z_{max}$ s krokem $dz$}
    \For{každý trojúhelník $(p_1, p_2, p_3) \in DT$}
        \State $\Delta z_1 \gets z - z_1$, $\Delta z_2 \gets z - z_2$, $\Delta z_3 \gets z - z_3$
        \If{všechna $\Delta z_i = 0$} \textbf{continue}
        \ElsIf{$\Delta z_1 \cdot \Delta z_2 \leq 0$ \textbf{and} $\Delta z_2 \cdot \Delta z_3 \leq 0$}
            \State vytvoř segment z hran $(p_1, p_2)$ a $(p_2, p_3)$
        \ElsIf{$\Delta z_2 \cdot \Delta z_3 \leq 0$ \textbf{and} $\Delta z_3 \cdot \Delta z_1 \leq 0$}
            \State vytvoř segment z hran $(p_2, p_3)$ a $(p_3, p_1)$
        \ElsIf{$\Delta z_3 \cdot \Delta z_1 \leq 0$ \textbf{and} $\Delta z_1 \cdot \Delta z_2 \leq 0$}
            \State vytvoř segment z hran $(p_3, p_1)$ a $(p_1, p_2)$
        \EndIf
    \EndFor
\EndFor
\end{algorithmic}
\end{algorithm}

\subsection{Analýza sklonu}
\begin{algorithm}[H] \caption{Analyze Slope}
\begin{algorithmic}[1]
\Require triangulace $DT$
\Ensure seznam trojúhelníků s přiřazeným sklonem
\For{každý trojúhelník $(p_1, p_2, p_3) \in DT$}
    \State spočti normálu $\vec{n} = (a, b, c) = \vec{p_1 p_2} \times \vec{p_1 p_3}$
    \State $|\vec{n}| \gets \sqrt{a^2 + b^2 + c^2}$
    \State $\varphi \gets \arccos(|c| / |\vec{n}|)$
    \State přiřaď trojúhelníku hodnotu $\varphi$
\EndFor
\end{algorithmic}
\end{algorithm}

\subsection{Analýza expozice}
\begin{algorithm}[H] \caption{Analyze Aspect}
\begin{algorithmic}[1]
\Require triangulace $DT$
\Ensure seznam trojúhelníků s přiřazenou expozicí
\For{každý trojúhelník $(p_1, p_2, p_3) \in DT$}
    \State spočti normálu $\vec{n} = (a, b, c) = \vec{p_1 p_2} \times \vec{p_1 p_3}$
    \State $A \gets \text{atan2}(a, -b)$ \Comment{negace $b$ kvůli převrácené Y ose na obrazovce}
    \If{$A < 0$} $A \gets A + 2\pi$ \EndIf
    \State přiřaď trojúhelníku hodnotu $A$
\EndFor
\end{algorithmic}
\end{algorithm}

\newpage
\section{Vstupní data a Aplikace}
Vstupními daty je textový soubor .txt obsahující trojice souřadnic bodů ve formátu $X, Y, Z$ oddělených bílými znaky. Pro testování byla použita data exportovaná z digitálního modelu reliéfu INSPIRE EL GRID s rozlišením $5 \times 5$ m. Tato data byla oříznuta podle masky oblasti Kladna (clip raster by mask), následně převedena z rastrového formátu na body (raster to points) a z výsledné sady byl vytvořen náhodný výběr (subset features) o velikosti 420 bodů. Minimální doporučený počet bodů dle zadání je 300.

Při načítání dat je proveden následující postup: souřadnice jsou načteny z textového souboru; je nalezen rozsah hodnot ($x_{min}, x_{max}, y_{min}, y_{max}$); souřadnice jsou lineárně přeškálovány na velikost kresební plochy s okrajem 50 pixelů. Osa $Y$ je zároveň převrácena, protože na obrazovce narůstá směrem dolů, zatímco v geografických souřadnicích směrem vzhůru (sever). Výšková souřadnice $z$ zůstává zachována v původních hodnotách a využívá se při výpočtu vrstevnic a výškové stupnice.

Po spuštění skriptu \textit{MainForm.py} se otevře okno aplikace. V záložce \textit{File} jsou umístěny funkce \textit{Open} pro načtení textového souboru s body a \textit{Exit} pro ukončení aplikace. V záložce \textit{Analysis} jsou funkce \textit{Create DT} pro vytvoření Delaunay triangulace, \textit{Create contour lines} pro konstrukci vrstevnic, \textit{Analyze slope} pro analýzu sklonu a \textit{Analyze exposition} pro analýzu expozice terénu. Záložka \textit{View} obsahuje přepínače viditelnosti jednotlivých vrstev -- \textit{DT}, \textit{Contour lines}, \textit{Slope} a \textit{Exposition}. Poslední záložka \textit{Clear} obsahuje funkce \textit{Clear results} (zachová body a smaže analýzy) a \textit{Clear all} (smaže veškerá data). Všechny funkce mají ikonu na nástrojové liště.

Pokud není načten vstupní soubor, lze body zadat ručně kliknutím myší do kresební plochy, přičemž hodnota $z$ je generována náhodně v intervalu $(200, 600)$. Po načtení souboru se kliknutí myší automaticky zablokuje, aby nedošlo k nežádoucímu přidávání bodů do načtené množiny.

\begin{figure}[H] 
\centering 
\includegraphics[width=\textwidth]{GUI_ukazka.png} 
\caption{\textit{Ukázka GUI aplikace pro tvorbu DMT}} 
\end{figure}

\section{Tvorba spustitelného souboru}
V souladu se zadáním byla vytvořena plně funkční a přeložená aplikace ve formě binárního souboru (\texttt{.exe}). Tento spustitelný soubor pro operační systém Windows integruje veškeré nezbytné knihovny a datové složky, což umožňuje jeho okamžité spuštění bez nutnosti instalace interpretu jazyka Python. K vytvoření binárního souboru byl využit nástroj \textit{Auto-py-to-exe}, což je grafické rozhraní nad knihovnou \textit{PyInstaller}. 

\section{Dokumentace}

\subsection{Třída \texttt{QPoint3DF} (\texttt{qpoint3df.py})}
Třída rozšiřuje \texttt{QPointF} z knihovny PyQt6 o třetí souřadnici $z$, což umožňuje reprezentaci bodů v 3D prostoru při zachování kompatibility s funkcemi pro vykreslování 2D.

\subsection{Třída \texttt{Edge} (\texttt{edge.py})}
Třída reprezentuje orientovanou hranu zadanou dvojicí bodů (počátečním a koncovým). Obsahuje metody pro získání bodů, přehození orientace a test rovnosti dvou hran.

\subsection{Třída \texttt{Triangle} (\texttt{triangle.py})}
Třída reprezentuje trojúhelník zadaný třemi vrcholy a dále uchovává jeho vypočtené atributy -- sklon (\texttt{slope}) a expozici (\texttt{aspect}). Slouží jako datová struktura pro výsledky analýz.

\subsection{Třída \texttt{Algorithms} (\texttt{algorithms.py})}
Třída obsahuje veškerou algoritmickou logiku úlohy.

\begin{itemize}
\item \texttt{getPointLinePosition(a, b, p)} -- určí polohu bodu vůči orientované přímce (test poloroviny pomocí determinantu).
\item \texttt{getNearestPoint(p, points)} -- nalezne nejbližší bod k zadanému bodu.
\item \texttt{get2LinesAngle(p1, p2, p3, p4)} -- spočítá úhel mezi dvěma přímkami pomocí skalárního součinu.
\item \texttt{findDelaunayPoint(p1, p2, points)} -- nalezne bod splňující Delaunayovo kritérium pro danou hranu.
\item \texttt{createDT(points)} -- vytvoří Delaunay triangulaci inkrementální metodou se seznamem aktivních hran.
\item \texttt{updateAEL(e, AEL)} -- aktualizuje seznam aktivních hran.
\item \texttt{getContourPoint(p1, p2, z)} -- vypočte průsečík hrany s horizontální rovinou $z$.
\item \texttt{createContourLines(DT, z\_min, z\_max, dz)} -- vytvoří vrstevnice v daném intervalu s daným krokem pomocí lineární interpolace.
\item \texttt{createContourLineSegment(p1, p2, p3, z, contour\_lines)} -- vytvoří segment vrstevnice v jednom trojúhelníku.
\item \texttt{getPlaneNormal(p1, p2, p3)} -- spočítá normálový vektor roviny procházející třemi body.
\item \texttt{analyzeSlope(DT)} -- provede analýzu sklonu všech trojúhelníků.
\item \texttt{analyzeAspect(DT)} -- provede analýzu expozice všech trojúhelníků.
\end{itemize}

\subsection{Třída \texttt{Draw} (\texttt{draw.py})}
Třída zajišťuje veškeré grafické operace -- dědí z \texttt{QWidget}.

\begin{itemize}
\item \texttt{mousePressEvent(e)} -- přidá bod kliknutím (pouze pokud není načten soubor).
\item \texttt{paintEvent(e)} -- vykreslí v tomto pořadí: vyplněné trojúhelníky sklonu, vyplněné trojúhelníky expozice, hrany DT, běžné vrstevnice tenčí světlejší hnědou čarou, hlavní (zvýrazněné) vrstevnice silnější tmavší hnědou čarou, body.
\item \texttt{getAspectColor(aspect)} -- vrací jednu z 8 pevně definovaných RGB barev odpovídající příslušné světové straně (N, SV, V, JV, J, JZ, Z, SZ).
\item \texttt{setPoints(points)} / \texttt{setDT(DT)} / \texttt{setContours(contours)} / \texttt{setMainContours(contours)} / \texttt{setSlopeTriangles(t)} / \texttt{setAspectTriangles(t)} -- settery pro jednotlivé datové struktury.
\item \texttt{setFileLoaded(value)} -- zablokuje/povolí přidávání bodů myší.
\item \texttt{setShowDT/Contours/Slope/Aspect/Points(value)} -- přepínače viditelnosti vrstev.
\item \texttt{clearResult()} -- smaže výsledky analýz, zachová body.
\item \texttt{clearAll()} -- smaže všechna data včetně bodů.
\end{itemize}

\subsection{Třída \texttt{Ui\_MainWindow} (\texttt{MainForm.py})}
Třída inicializuje hlavní okno aplikace a propojuje třídy \texttt{Draw} a \texttt{Algorithms}.

\begin{itemize}
\item \texttt{createDTClick()} -- vytvoří Delaunay triangulaci nad načtenými body.
\item \texttt{createContourLinesClick()} -- vytvoří vrstevnice, rozsah $z_{min}, z_{max}$ se automaticky určí z dat, krok $dz$ se volí tak, aby vzniklo přibližně 20 vrstevnic.
\item \texttt{analyzeSlopeClick()} -- provede analýzu sklonu a vizualizuje výsledek.
\item \texttt{analyzeExpositionClick()} -- provede analýzu expozice a vizualizuje výsledek.
\item \texttt{openFileClick()} -- otevře dialog pro výběr souboru, načte body, přeškáluje souřadnice na canvas a zablokuje přidávání bodů myší.
\item \texttt{clearResultClick()} / \texttt{clearAllClick()} -- mazací funkce.
\item \texttt{toggleDT()} / \texttt{toggleContours()} / \texttt{toggleSlope()} / \texttt{toggleExposition()} -- přepínače viditelnosti jednotlivých vrstev.
\end{itemize}

\newpage
\section{Výsledky}
Testování aplikace proběhlo nad datovou sadou obsahující 420 geodetických bodů oblasti Kladna získaných z modelu INSPIRE EL GRID s rozlišením $5 \times 5$\,m (rozsah $z$ přibližně 275--450\,m n.\,m.).

\begin{figure}[H] 
\centering 
\includegraphics[width=0.85\textwidth]{DT_ukazka.png} 
\caption{\textit{Delaunay triangulace}} 
\end{figure}

\begin{figure}[H] 
\centering 
\includegraphics[width=0.85\textwidth]{vrstevnice.png} 
\caption{\textit{Vrstevnice vytvořené lineární interpolací v triangulaci}} 
\end{figure}

\begin{figure}[H] 
\centering 
\includegraphics[width=0.85\textwidth]{sklon_ukazka.png} 
\caption{\textit{Vizualizace analýzy sklonu}} 
\end{figure}

\begin{figure}[H] 
\centering 
\includegraphics[width=0.85\textwidth]{expozice_ukazka.png} 
\caption{\textit{Vizualizace analýzy expozice}} 
\end{figure}

Delaunay triangulace byla úspěšně vytvořena nad datovou sadou. Konstrukce vrstevnic s automaticky určeným intervalem (přibližně 20 vrstevnic napříč rozsahem $z$) poskytuje dobře čitelnou vizualizaci reliéfu. Analýza sklonu zřetelně odhaluje strmé oblasti v kopcovitém terénu, analýza expozice pak přehledně ukazuje orientaci jednotlivých svahů.

\section{Zhodnocení a slabiny algoritmu}
Algoritmus založený na 2D Delaunay triangulaci poskytuje kvalitní model terénu v převážné většině případů. Existují však situace, kdy jeho výsledky nejsou optimální nebo selhávají.

\textbf{Plochy se svislou nebo převislou složkou} -- 2D Delaunay triangulace funguje v rovině $xy$ a předpokládá, že každému bodu odpovídá právě jedna výška $z$. Terénní útvary jako jsou převisy, jeskyně nebo svislé skalní stěny tak nelze korektně reprezentovat. Model v tomto případě vytvoří plochu, která neodpovídá realitě.

\textbf{Nekonvexní terény a dlouhé úzké údolí} -- Algoritmus standardně vyplní celou konvexní obálku vstupní množiny. U nekonvexních území (ostrovní území, údolí, meandrující řeky) tak vznikají trojúhelníky přes oblasti, kde body ve skutečnosti nejsou, což vede k nesprávné interpolaci výšek.

\textbf{Ostré hřbety a údolnice} -- Delaunay triangulace nezná pojem charakteristické čáry terénu (hřbet, údolnice, hrana). Pokud nejsou tyto linie vyznačeny vstupními body dostatečně hustě, algoritmus je neidentifikuje a vyhladí, což vede ke ztrátě důležitých terénních rysů.

\textbf{Nerovnoměrně rozmístěné body} -- V oblastech s řídkým rozmístěním bodů mohou vznikat velké trojúhelníky, které hrubě aproximují terén. Sklon a expozice těchto rozsáhlých ploch pak nemusí odpovídat skutečnosti.

\textbf{Kolineární body} -- Pokud leží více bodů na jedné přímce, může docházet k numerickým nestabilitám při výběru Delaunay bodu. V naší implementaci je tento případ ošetřen tolerancí v testu polohy.

\section{Závěr}
V rámci zpracování úlohy byla implementována Delaunay triangulace metodou inkrementální konstrukce se seznamem aktivních hran, konstrukce vrstevnic pomocí lineární interpolace ve zvoleném intervalu a s automaticky určeným krokem, analýza sklonu terénu a analýza expozice terénu. Grafické rozhraní bylo realizováno v PyQt6 a umožňuje interaktivní práci s daty -- jejich načtení z textového souboru, postupné spouštění jednotlivých analýz, přepínání viditelnosti vrstev a mazání výsledků.

Aplikace byla úspěšně otestována nad geodetickými daty INSPIRE EL GRID pro oblast Kladna. Pro vizualizaci expozice byla zvolena diskrétní barevná paleta s 8 barvami odpovídajícími hlavním světovým stranám: sever světle zelená, severovýchod tmavě zelená, východ modrá, jihovýchod magenta, jih růžová, jihozápad červená, západ oranžová a severozápad žlutá. Pro vizualizaci sklonu byla použita stupnice šedi, kde světlá barva odpovídá ploché ploše a tmavá strmému svahu. Jako možné budoucí vylepšení se nabízí implementace triangulace nekonvexní oblasti (ořezání hranicí polygonu), automatický popis vrstevnic respektující kartografické zásady, barevná hypsometrie a 3D vizualizace terénu s využitím promítání.

\newpage
\renewcommand{\refname}{Zdroje} 
\begin{thebibliography}{9}

\bibitem{bayer2026} BAYER, T. (2026): \textit{Digitální model terénu}. Přednáška pro předmět Algoritmy počítačové kartografie, Katedra aplikované geoinformatiky a kartografie, Přírodovědecká fakulta UK.
\bibitem{inspire} Geoportál INSPIRE (2026): \textit{Výšková data INSPIRE EL GRID}. Dostupné z: \url{https://geoportal.gov.cz/}

\end{thebibliography}
\end{document}